\addcontentsline{toc}{section}{ВВЕДЕНИЕ}
\begin{center}
    \textbf{ВВЕДЕНИЕ}
        \vspace{-1em}
\end{center}

В последние годы наблюдается стремительный рост объёмов данных, представляемых в виде графов. Социальные сети, системы рекомендаций, телекоммуникационные сети, транспортные системы, базы знаний и многие другие современные информационные системы естественным образом описываются графовыми структурами. В отличие от традиционных реляционных моделей данных, графы позволяют эффективно представлять сложные взаимосвязи между объектами и выполнять запросы, ориентированные на поиск путей, анализ связности, обнаружение сообществ и выявление закономерностей в структуре данных.

Рост масштабов графовых данных приводит к необходимости их хранения и обработки в распределённых вычислительных системах. Современные графовые базы данных способны содержать миллиарды вершин и рёбер, что делает невозможным размещение всего графа на одном вычислительном узле. В таких условиях данные распределяются между несколькими хранилищами, образующими единое логическое пространство хранения. Однако переход к распределённой архитектуре порождает новую проблему — необходимость эффективного распределения вершин графа между узлами хранения.

Качество распределения графовых данных оказывает непосредственное влияние на производительность распределённой системы. При выполнении запросов поиска пути, обхода графа или анализа связей система вынуждена обращаться к данным, находящимся на различных узлах. Если вершины, участвующие в одном запросе, располагаются на разных шардах, возникает необходимость передачи данных по сети и синхронизации между вычислительными узлами. Такие межшардовые обращения являются существенно более дорогостоящими по сравнению с локальными операциями и часто становятся основным фактором, ограничивающим производительность распределённой графовой базы данных.

Задача распределения графовых данных представляет собой сложную оптимизационную проблему. С одной стороны, необходимо минимизировать количество рёбер, соединяющих вершины различных шардов, поскольку именно они приводят к межузловым взаимодействиям. С другой стороны, требуется обеспечить равномерное распределение нагрузки между узлами хранения, не допуская существенного дисбаланса размеров шардов. Дополнительную сложность создаёт динамический характер современных информационных систем, в которых граф непрерывно изменяется вследствие добавления новых вершин и рёбер.

Для решения данной задачи было предложено множество подходов. Классические методы структурного разбиения графов, такие как METIS, ориентированы на минимизацию разреза графа и используют его топологические свойства для формирования качественного распределения. Более современные методы учитывают не только структуру графа, но и особенности рабочей нагрузки, анализируя реальные запросы пользователей и стремясь минимизировать стоимость их выполнения. К таким подходам относятся алгоритмы TAPER, Loom, Schism, SWORD и WAWPart. Отдельное направление исследований составляют методы потокового распределения, предназначенные для размещения новых элементов в уже распределённом графе без необходимости полного перераспределения данных.

Несмотря на значительное количество существующих решений, проблема эффективного распределения графовых данных остаётся актуальной. Многие алгоритмы обладают высокой вычислительной сложностью, требуют полного анализа графа или плохо адаптируются к динамическому изменению структуры данных. Кроме того, часть существующих методов ориентирована на пакетную обработку графов и не предназначена для работы в условиях непрерывного поступления новых данных. В связи с этим возникает необходимость разработки программных средств, позволяющих эффективно распределять вершины графа между хранилищами с учётом особенностей современных распределённых графовых баз данных.

Актуальность темы выпускной квалификационной работы обусловлена широким распространением графовых баз данных и необходимостью повышения эффективности их функционирования в распределённой среде. Снижение количества межшардовых переходов позволяет уменьшить сетевые накладные расходы, сократить время выполнения запросов и повысить масштабируемость системы в целом. Разработка специализированного программного модуля распределения вершин представляет практический интерес как для научных исследований в области распределённых вычислений, так и для создания высокопроизводительных графовых информационных систем.

Объектом исследования являются распределённые графовые базы данных.

Предметом исследования являются алгоритмы и методы распределения вершин графа между узлами хранения в распределённой системе.

Целью работы является разработка программного модуля распределения вершин в графовой базе данных, обеспечивающего уменьшение количества межшардовых переходов при сохранении балансировки распределения данных между хранилищами.

Для достижения поставленной цели необходимо решить следующие задачи:

\begin{enumerate}
\item выполнить анализ существующих методов распределения графовых данных и алгоритмов разбиения графов;
\item исследовать методы структурной оптимизации распределения, методы распределения с учётом рабочей нагрузки и методы потокового распределения;
\item выбрать алгоритмы, наиболее подходящие для решения поставленной задачи;
\item разработать архитектуру программного модуля распределения вершин графа;
\item реализовать программный модуль и необходимые структуры данных;
\item разработать механизм имитации распределённой графовой системы для проведения вычислительных экспериментов;
\item провести тестирование разработанного программного обеспечения и оценить эффективность предложенного решения.
\end{enumerate}

В рамках работы рассматривается задача распределения вершин графа по множеству хранилищ таким образом, чтобы минимизировать количество межшардовых переходов при выполнении запросов поиска пути. В качестве критерия качества распределения используется число рёбер найденных путей, соединяющих вершины, расположенные на различных шардах. Данная метрика напрямую характеризует объём межузлового взаимодействия и позволяет оценивать эффективность алгоритмов независимо от особенностей реализации конкретной графовой базы данных.

Научная новизна работы заключается в разработке программного модуля, объединяющего методы начального распределения графовых данных и механизмы последующего улучшения распределения при изменении структуры графа. Предлагаемый подход ориентирован на использование в распределённых графовых системах и учитывает необходимость сохранения баланса между качеством распределения и вычислительными затратами на его построение.

Практическая значимость работы состоит в возможности использования разработанного программного модуля для исследования и сравнения различных алгоритмов распределения графовых данных, а также для повышения эффективности функционирования распределённых графовых баз данных за счёт уменьшения количества межшардовых переходов.

В первом разделе работы выполнен анализ существующих методов распределения графовых данных, рассмотрены классические и современные подходы к разбиению графов, исследованы методы структурной оптимизации, методы, учитывающие рабочую нагрузку, а также алгоритмы потокового распределения. Во втором разделе представлены проектирование архитектуры программного модуля, выбор используемых алгоритмов и особенности их программной реализации. Третий раздел посвящён разработке методики тестирования, проведению вычислительных экспериментов и анализу полученных результатов. В заключении приведены основные результаты работы и направления дальнейшего развития разработанного программного решения.
